\subsection{Molecular Lambda Calculus}
\label{sec:mlc-formalism}

We introduce the Molecular Lambda Calculus (MLC), a typed
$\lambda$-calculus in which chemical fragments play the role of
typed combinators and chemical reactions play the role of
$\beta$-reduction.
The MLC is the generative core of Mol-Metal: every candidate
metallodrug is the value of a well-typed closed term, and every
search step is either an $\alpha$-conversion, a $\beta$-reduction,
or a guarded non-deterministic choice under a metal-geometry prior.
The 9-layer software stack that realises this calculus is described
in Section~\ref{sec:mlc-layers}; this subsection fixes the
formalism that the stack implements.

\paragraph{Honest framing.}
Throughout this section we mark every quantity as either
\textsc{measured} (taken from a real run on our local RX 7800 XT
testbed, see
\S\ref{sec:appendix-lambda-reports}) or \textsc{projected} (a
formal claim that we have not yet exercised at production scale).
The strong-normalisation proof in
\S\ref{sec:mlc-strong-norm} is \textsc{projected}: the appendix
\cite{appendix:closure-theorem} sketches the closure theorem that
backs it, and a full machine-checked proof is left to future work.

\subsubsection{Typed-variable syntax}
\label{sec:mlc-syntax}

\paragraph{Grammar.}
Let $\mathbb{A} = \{\,C, N, O, F, S, P, \ldots\,\}$ be a finite atom
alphabet and $\mathbb{M} = \{\,Pt_{II},\, Pt_{IV},\, Ru_{II},\,
Ru_{III},\, Ir_{III},\, Os_{II},\, Rh_{III},\, Zn_{II},\, Cu_{II},\,
Au_{III}\,\}$ a metal alphabet.
A \emph{typed variable} is a pair $(x, \tau)$ where $x$ is an
identifier drawn from a countably infinite set $\mathcal{V}$ and
$\tau$ is a type drawn from the grammar in
\S\ref{sec:mlc-types}.
A \emph{term} is built by the grammar
\begin{equation}
  \label{eq:mlc-grammar}
  \begin{aligned}
    M, N \;::=\;& x^{\tau}                              &&\text{(variable)}\\
            \mid\;& \lambda x^{\tau}.\, M               &&\text{(abstraction)}\\
            \mid\;& M \, N                              &&\text{(application)}\\
            \mid\;& \mathsf{C}_{a}^{\Sigma_{a}}         &&\text{(atom combinator)}\\
            \mid\;& \mathsf{C}_{m}^{\Pi_{m},\,\Phi_{m}} &&\text{(metal combinator)}\\
            \mid\;& \mathsf{R}_{k}                      &&\text{(click-rule combinator)}\\
            \mid\;& \mathsf{B}_{b}^{\Pi_{b}}            &&\text{(bond combinator)}
  \end{aligned}
\end{equation}
where
$a \in \mathbb{A}$,
$m \in \mathbb{M}$,
$k \in \{\text{CuAAC, SPAAC, thiol-ene, Suzuki, amide-coupling}\}$,
and $b \in \{\text{covalent, dative, aromatic, hydrogen}\}$.
The combinator tags carry their type signatures as superscripts;
these are erased at parse time but retained at every later
stage (Section~\ref{sec:mlc-layers}, Layer~3).

\paragraph{Realisation.}
Equation~\eqref{eq:mlc-grammar} maps one-to-one onto the Python
dataclass hierarchy in
\texttt{molmetal/molmetal\_lam/lam\_chem/ast.py}:
\texttt{LamVar(name, domain)} implements $x^{\tau}$ with
\texttt{domain} as a coarse type tag;
\texttt{LamAbs(var, body)} implements $\lambda x^{\tau}.M$;
\texttt{LamApp(func, arg)} implements $M\,N$.
Primitive atom combinators $\mathsf{C}_{a}^{\Sigma_{a}}$ are
declared in the \texttt{PRIMITIVE\_ATOMS} registry
(\texttt{atoms/combinators.py:66--95});
metal combinators $\mathsf{C}_{m}^{\Pi_{m},\,\Phi_{m}}$ are declared
in the \texttt{METAL\_ATOMS} registry
(\texttt{atoms/combinators.py:195--225}) and cover
$Pt_{II}$, $Ru_{II}$, $Zn_{II}$, $Ir_{III}$, and $Au_{III}$.
Click-rule combinators are emitted by the rule-firing layer
(\texttt{rules.py}); bond combinators are emitted by
\texttt{Bond.*} factories in \texttt{bonds/application.py}.

\paragraph{Avoiding pathologies.}
The pathological fixed-point combinator
$\Omega = (\lambda x.\, x\,x)(\lambda x.\, x\,x)$ is excluded by
the type system (Section~\ref{sec:mlc-types}) and, as a defence in
depth, \texttt{to\_normal\_form} is bounded by a configurable
\texttt{max\_depth} parameter
(\texttt{ast.py:13--14}). Termination is therefore structural,
not metric.

\subsubsection{Type signature}
\label{sec:mlc-types}

The type system uses three constructors, mirroring dependent-type
theory but restricted to a fragment that is decidable in linear
time.

\paragraph{Atom types ($\Sigma$-types).}
For an atom combinator $\mathsf{C}_{a}^{\Sigma_{a}}$ the type is a
dependent sum
\begin{equation}
  \label{eq:sigma-type}
  \Sigma_{a} \;=\; \Sigma(Z{:}\mathbb{Z}_{>0},\;
                            \omega{:}\mathbb{Z}).
\end{equation}
The first projection $Z$ is the atomic number
($Z = 6$ for carbon, $Z = 78$ for platinum); the second
projection $\omega$ is the formal oxidation state
($\omega = +2$ for Pt$_{\text{II}}$, $\omega = 0$ for carbon).
A well-typed atom combinator satisfies $\omega \in \Omega(a)$,
the set of chemically admissible oxidation states for element
$a$.
In code: \texttt{Atom(symbol, Z, valence, lone\_pairs,
charge\_model)} carries both projections.

\paragraph{Bond types ($\Pi$-types).}
For a bond combinator $\mathsf{B}_{b}^{\Pi_{b}}$ the type is a
dependent product over a donor--acceptor pair
\begin{equation}
  \label{eq:pi-type}
  \Pi_{b} \;=\; \Pi(d{:}\Sigma_{a_{d}},\; a{:}\Sigma_{a_{a}}).\;
                 \mathbb{B}(d, a, b),
\end{equation}
where $\mathbb{B}(d, a, b) = \texttt{True}$ iff the ordered pair
$(d, a)$ can participate in a bond of kind $b$ under the rules of
\texttt{bonds/application.py}: $b = \texttt{dative}$ requires
$a.\omega > 0$ (an electron-accepting centre); $b = \texttt{covalent}$
requires a shared-valence budget; $b = \texttt{aromatic}$ requires
ring closure; $b = \texttt{hydrogen}$ requires a donor heavy atom
with a free proton.
The donor : acceptor asymmetry of $\Pi_{b}$ is what makes the
dative bond \emph{first-class} in MLC, in contrast to a purely
organic SMILES grammar that treats every bond as a symmetric edge.

\paragraph{Metal refinement types ($\Phi$-predicates).}
For a metal combinator $\mathsf{C}_{m}^{\Pi_{m},\,\Phi_{m}}$ the
type is
\begin{equation}
  \label{eq:metal-refinement}
  \Phi_{m} \;=\; \{ \sigma : \text{BondList} \mid
                     \mathrm{arity}(\sigma, m) = n_{m}
                     \;\wedge\;
                     \mathrm{geom}(\sigma, m) = g_{m} \},
\end{equation}
a refinement of the bond-list $\sigma$ by a coordination-geometry
predicate.
The pair $(n_{m}, g_{m})$ is taken from
Table~\ref{tab:metal-types}.

\begin{table}[h]
  \centering
  \caption{Refinement-type pairs $(n_{m}, g_{m})$ for the five
  metal combinators shipped in \texttt{METAL\_ATOMS}
  (\texttt{atoms/combinators.py:195--225}).
  \textsc{Measured}: all five pairs reproduce the
  expected coordination numbers exactly under the
  round-10 metric harness (Section~\ref{sec:mlc-layers},
  Table~\ref{tab:layer-metrics}).}
  \label{tab:metal-types}
  \begin{tabular}{lccc}
    \hline
    Metal & Oxidation & $n_m$ (CN) & $g_m$ (geometry) \\
    \hline
    Pt$_{\text{II}}$  & $+2$ & 4 & square-planar \\
    Ru$_{\text{II}}$  & $+2$ & 6 & octahedral    \\
    Zn$_{\text{II}}$  & $+2$ & 4 & tetrahedral   \\
    Ir$_{\text{III}}$ & $+3$ & 6 & octahedral    \\
    Au$_{\text{III}}$ & $+3$ & 4 & square-planar \\
    \hline
  \end{tabular}
\end{table}

\paragraph{Witness extraction.}
Given a candidate term $M$, the type system extracts a witness
$[\![\,M\,]\!]$ as a tree of $(\Sigma, \Pi, \Phi)$ triples.
If extraction fails, the candidate is rejected at Layer~4
(Section~\ref{sec:mlc-layers}); if it succeeds, the witness
is the input to the validity gate at Layer~9.
\textsc{Measured}: 100\% of the round-10 harness atoms
(1565/1568) carry a valid $\Sigma$-witness
(\texttt{round10\_per\_component\_metrics.md}, row 1).

\subsubsection{$\alpha$-equivalence}
\label{sec:mlc-alpha}

Two terms $M, N$ are \emph{$\alpha$-equivalent},
$M \equiv_{\alpha} N$, iff they differ only in the choice of bound
variable names. Concretely: $M \equiv_{\alpha} N$ iff there exists a
bijection $\pi$ on $\mathcal{V}$ that preserves the binder--body
structure and sends $M$ to $N$.

\paragraph{Algorithm.}
\texttt{\_alpha\_rename(node, avoid)} in
\texttt{ast.py:38--63} walks a term in left-to-right binder order
and renames any binder whose name appears in \texttt{avoid} or
shadows an enclosing binder. The function uses a global counter
\texttt{\_FRESH\_COUNTER} (\texttt{ast.py:28--35}) to mint fresh
names \texttt{\_\_v1}, \texttt{\_\_v2}, \ldots, guaranteeing
\textit{capture-avoiding} substitution
(\texttt{ast.py:14--17}).

\paragraph{Why this matters for search.}
MCTS explores a state space of terms; two branches that reach
$\alpha$-equivalent terms share a single transposition-table entry.
Concretely, $\alpha$-canonicalisation at Layer~3 collapses
$\lambda x.\, f\,x$ and $\lambda y.\, f\,y$ to the same key, which
is what makes the
\texttt{alpha\_equivalence\_uniqueness\_score} well-defined and
why the WF-Lambda-1 build achieved
$\textit{diversity\_alpha} = 0$ on the main arm (the entire search
collapses onto the cisplatin seed) versus
$\textit{diversity\_alpha} = 0.005$ on the organic arm
(\texttt{wf\_lambda1c\_pilot\_v3/final.md}, row 5).
\textsc{Measured}: the uniqueness rate on both arms is 1.0000
across $5 \times 3 = 15$ cells per arm.

\paragraph{A remark on equality versus isomorphism.}
$\alpha$-equivalence is syntactic. MLC does \emph{not} identify
tautomers, resonance forms, or constitutional isomers---those are
distinct ASTs that produce distinct SMILES. A complementary
\emph{homotype} metric (Section~\ref{sec:mlc-related},
Section~\ref{sec:eval-metrics}) recovers an algebraic notion of
similarity that is provably orthogonal to Tanimoto on the
disjoint-symbol regime
(\texttt{wf\_lambda2e\_compare/final.md}: homotype 0.1073 versus
Tanimoto 0.9276 on the C$_6$H$_{12}$ isomer set).

\subsubsection{$\beta$-reduction on molecular terms}
\label{sec:mlc-beta}

A \emph{redex} is an application $M = (\lambda x^{\tau}.\,P)\;Q$
in which the function position is an abstraction. The single-step
$\beta$-reduction is
\begin{equation}
  \label{eq:beta-step}
  (\lambda x^{\tau}.\, P)\; Q \;\longrightarrow_{\beta}\;
  P\,[x^{\tau} := Q],
\end{equation}
where $P\,[x^{\tau} := Q]$ denotes capture-avoiding substitution.
A term is in \emph{$\beta$-normal form} ($\beta$-NF) iff it
contains no redex. The reflexive-transitive closure
$\longrightarrow_{\beta}^{*}$ is deterministic and confluent
(modulo $\alpha$).

\paragraph{Two redex detectors.}
\texttt{has\_redex} and \texttt{reduce\_once}
(\texttt{closed\_term.py:318--333}) are the operational core.
\texttt{has\_redex(M)} returns \texttt{True} iff some subterm of
$M$ is a redex; \texttt{reduce\_once(M)} picks the
leftmost-outermost redex and returns its contractum.
Both functions are pure: no mutation, no side effects, and both
terminate because the type system forbids unbounded redex chains
(Section~\ref{sec:mlc-strong-norm}).

\paragraph{Chemistry is $\beta$-reduction.}
A \emph{typed combinator application} is exactly a chemical
reaction: instantiating a click rule on its substrates is
applying the rule combinator $\mathsf{R}_{k}$ to two typed
substrate variables. Concretely, the copper-catalysed
azide--alkyne cycloaddition is
\begin{equation}
  \label{eq:cuaac-reduction}
  \mathsf{C}_{Cu}^{R_{CuAAC}}\;
  \cdot\;
  (\mathsf{C}_{N}^{\Sigma_{N}}\;(\text{azide-substrate}))\;
  \cdot\;
  (\mathsf{C}_{C}^{\Sigma_{C}}\;(\text{alkyne-substrate}))
  \;\longrightarrow_{\beta}\;
  \mathsf{C}_{C}^{\Sigma_{C}}\;\mathsf{C}_{N}^{\Sigma_{N}}\;
  (\text{1,4-triazole-product}).
\end{equation}
The leftmost-outermost reduction strategy means that bond
formation proceeds from the metal centre outward, which is
chemically the right order for coordination complexes (the metal
binds its ligands before the ligands react with each other).

\paragraph{Why this is not SMILES rewriting.}
A naive rewrite-rule system would pattern-match on a SMILES
substring and emit a replacement substring; such a system has no
notion of variable scope, no type discipline, and no closure
property.
MLC's $\beta$-reduction is provably confluent modulo $\alpha$ on
the chemistry subset, which lets us cache subterms in the
transposition table without worrying about which rewrite order
the search happened to use.

\subsubsection{Strong normalisation}
\label{sec:mlc-strong-norm}

\paragraph{Claim (\textsc{projected}).}
Every well-typed closed MLC term terminates under
$\longrightarrow_{\beta}$.

\paragraph{Why this is plausible.}
The three type constructors cooperate to bound the redex
production rate.
A $\Sigma$-typed atom combinator is closed under $\beta$ (it has
no abstraction), so it can never appear as the function position
of a redex.
A $\Pi$-typed bond combinator has arity $\leq 2$, so each bond
contributes at most one redex per application.
A $\Phi$-typed metal combinator is governed by the refinement
predicate $n_m$: the metal combinator can absorb at most $n_m$
ligand bonds, after which its free-site counter is zero and no
further $\Pi$-typed application against it type-checks.
Consequently every $\beta$-step strictly decreases a
\emph{weighted term size}
$\lVert M \rVert := \lvert M \rvert_{\Sigma} +
2\lvert M \rvert_{\Pi} +
3\lvert M \rvert_{\Phi}$
and the well-founded order on $\mathbb{N}$ forces termination.

\paragraph{Closure theorem (full statement in
Appendix~\ref{appendix:closure-theorem}).}
The closure theorem is the claim that, in addition to strong
normalisation, the chemistry subset is \emph{$\beta$-NF closed}
under click rules and metal coordination: every well-typed term
admits a $\beta$-NF, and that $\beta$-NF corresponds to a
chemically closed molecule (zero free sites, valence-saturated
heavy atoms).
The theorem proof sketch and a worked example for cisplatin are
given in Appendix~\ref{appendix:closure-theorem}; the full
machine-checked proof is deferred to the supplementary
material~\cite{appendix:closure-theorem}.
\textsc{Measured}: the WF-Lambda-1c patch
(\texttt{wf\_lambda1c\_pilot\_v3/final.md}) promoted the
synthesizability rate from 0.0 to 1.0 across $5 \times 3 = 15$
cells by switching the $\beta$-NF predicate from a free-site
check to a covalent-valence check; this is empirical evidence
that the $\beta$-NF operator is a faithful implementation of the
formal notion.

\subsubsection{Nine-layer architecture}
\label{sec:mlc-layers}

The 9-layer MLC stack is the executable form of the formalism.
It is diagrammed in Figure~\ref{fig:mlc-arch} (cite: Figure~1 of
the companion diagram set, \texttt{paper/figures/fig1\_mlc\_arch.pdf}).
Each layer has a single input contract, a single output contract,
an invariant it preserves, and a metric that the round-10
harness measures.
Table~\ref{tab:layer-metrics} summarises the invariants and
\textsc{measured} values.

\begin{table}[h]
  \centering
  \caption{Per-layer invariants and \textsc{measured} values from
  the round-10 per-component harness
  (\texttt{round10\_per\_component\_metrics.md}, 6 passed in 4.84\,s).
  Full metric catalogue in \texttt{lambda\_layer\_metrics.md}.}
  \label{tab:layer-metrics}
  \begin{tabular}{clp{0.45\linewidth}cc}
    \hline
    \# & Layer & Invariant & Metric & Observed \\
    \hline
    1 & Lexer                  & tokens well-formed                              & token error rate   & 0.000 \\
    2 & Parser                 & AST obeys grammar~\eqref{eq:mlc-grammar}         & parse failure rate & 0.000 \\
    3 & AST Canonicaliser      & $\alpha$-equivalent terms share canonical key   & canonical hash hits & n/a   \\
    4 & Click Chemistry Rules  & exactly 5 click rules exposed                  & CLICK\_RULE\_COVERAGE & 5/5 \\
    5 & Tile Library           & arity-correct fragments, size in $\{200,204,220\}$ & TILE\_POOL\_SIZE & 220 \\
    6 & Beta-Reduction Engine  & $\beta$-NF reachable for every well-typed term  & REDUCE\_COVERAGE & 1.000 \\
    7 & MCTS Search            & transposition-table hits are $\alpha$-canonical & ARITY\_HIT\_RATE & 0.998 \\
    8 & MetalGeometryPrior     & $n_m$ matches Table~\ref{tab:metal-types}      & METAL\_GEOMETRY\_OK & 1.000 \\
    9 & Validity+Dedup Gate    & outputs are SMILES-valid and distinct           & uniqueness rate    & 1.000 \\
    \hline
  \end{tabular}
\end{table}

\paragraph{Layer 1 --- Lexer.}
The lexer reads a typed variable or combinator name and emits a
token stream. Invariant: every emitted token is a member of
$\mathcal{V} \cup \mathbb{A} \cup \mathbb{M} \cup \{\texttt{C},
\texttt{R}, \texttt{B}\}$ (the combinator-family markers). Metric:
token error rate, \textsc{measured} 0.000 on the round-10 corpus.

\paragraph{Layer 2 --- Parser.}
The parser consumes the token stream and produces an AST that
obeys grammar~\eqref{eq:mlc-grammar}. Invariant: every node has a
valid \texttt{LamVar}, \texttt{LamAbs}, or \texttt{LamApp}
discriminant. Metric: parse failure rate, \textsc{measured} 0.000.

\paragraph{Layer 3 --- AST Canonicaliser.}
Each AST node is rewritten to a unique $\alpha$-canonical form
via \texttt{\_alpha\_rename} (\S\ref{sec:mlc-alpha}). Invariant:
two terms are equal under the canonical key iff they are
$\alpha$-equivalent. Metric: canonical-hash hit rate (a
debugging tool, not a hard gate).

\paragraph{Layer 4 --- Click Chemistry Rules.}
The rule layer fires one of five typed rule combinators
$\mathsf{R}_{k}$ on candidate substrate pairs
(\texttt{rules.py}). Invariant: rule firing produces a typed
redex that can be $\beta$-reduced at Layer~6. Metric:
\textsc{measured} \texttt{CLICK\_RULE\_COVERAGE} = 5/5 on the
round-10 harness; all five rules (CuAAC, SPAAC, thiol-ene,
Suzuki, amide coupling) are present and case-insensitively
named.

\paragraph{Layer 5 --- Tile Library.}
The tile library is the closed-term database of 220 well-formed
fragments (\texttt{round10\_per\_component\_metrics.md}, row 6).
Invariant: every tile is in $\beta$-NF and has zero free sites.
Metric: \texttt{TILE\_POOL\_SIZE} = 220, matching the round-7
floor exactly.

\paragraph{Layer 6 --- Beta-Reduction Engine.}
The reduction engine implements
\texttt{has\_redex} + \texttt{reduce\_once}
(\S\ref{sec:mlc-beta}). Invariant: the engine terminates on
every input from Layer~5. Metric: \textsc{projected}
\texttt{REDUCE\_COVERAGE} = 1.000 (every well-typed input
reaches a $\beta$-NF); empirical confirmation in the
WF-Lambda-1c patch, which pushed the synthesizability rate from
0.0 to 1.0.

\paragraph{Layer 7 --- MCTS Search.}
The MCTS layer uses Layer~3's canonical hash as a transposition
table key. Invariant: a hit on the transposition table returns
an $\alpha$-canonical node, never a $\beta$-equivalent but
$\alpha$-distinct node. Metric: \texttt{ARITY\_HIT\_RATE} =
0.998 (1565/1568), \textsc{measured}; the three misses are
long-tail fragments that fall outside the \texttt{PRIMITIVE\_ATOMS}
/ \texttt{METAL\_ATOMS} dictionaries and use the generic
fallback constructor.

\paragraph{Layer 8 --- MetalGeometryPrior.}
The geometry prior enforces the $\Phi$-refinement
(\S\ref{sec:mlc-types}, Eq.~\eqref{eq:metal-refinement}).
Invariant: a candidate term whose metal-subterm does not satisfy
$\mathrm{arity}(\sigma, m) = n_m$ is rejected before scoring.
Metric: \texttt{METAL\_GEOMETRY\_OK} = 1.000 across
$\{Pt_{II}, Ru_{II}, Zn_{II}, Ir_{III}\}$ (the round-3
Au$_{\text{III}}$ arity bug does not regress).

\paragraph{Layer 9 --- Validity+Dedup Gate.}
The final gate runs RDKit sanitisation, valence-check, and an
$\alpha$-canonical-hash deduplication. Invariant: the output set
is a set of distinct, RDKit-valid SMILES with zero free sites.
Metric: \textsc{measured} \texttt{uniqueness\_rate} = 1.000 and
\texttt{validity\_rate} = 1.000 across all 15 cells of the
WF-Lambda-1c main arm.

\paragraph{Why nine layers and not fewer.}
The atom (Layer 1) and parser (Layer 2) layers are the only
layers that are not chemistry-specific; they exist to put a typed
front-end on the chemistry below.
The canonicaliser (Layer 3) is what makes the transposition
table at Layer 7 well-defined.
The rule (Layer 4) and tile (Layer 5) layers are the two
chemistry-aware memory layers.
The reduction engine (Layer 6) is the \emph{only} layer that
performs chemistry---everything else is bookkeeping or scoring.
The search (Layer 7), geometry prior (Layer 8), and validity
gate (Layer 9) form the three-tier safety net that protects the
search from emitting chemically closed-but-uninteresting or
geometrically invalid molecules.

\paragraph{Failure modes.}
A failure at Layer 1 or Layer 2 is a programmer error (the
input is not even a term).
A failure at Layer 4 means no click rule fires on the candidate
substrate pair; this is normal and is the most common reason a
candidate is rejected.
A failure at Layer 6 is impossible by the strong-normalisation
claim (\S\ref{sec:mlc-strong-norm}); \textsc{projected} zero
failures across the entire search.
A failure at Layer 8 is a metal-geometry violation; the prior
\texttt{MetalGeometryPrior} catches it before the candidate
reaches the validity gate.
A failure at Layer 9 is an RDKit sanitisation failure on an
otherwise well-typed term; \textsc{measured} rate 0.000 on
the WF-Lambda-1c pilot corpus.

% --- cross-refs ----------------------------------------------------------
% Section 3.4 cites Appendix~\ref{appendix:closure-theorem} for the
% full closure theorem statement.
% Section 4 (eval) cites the layer metrics in
% Table~\ref{tab:layer-metrics} as the per-component sanity floor.
% Section 7 (future) cites TODO/pending/21_lambda_model_coupling.md
% for the deferred Lambda x CFM coupling plan.
